\documentclass[11pt]{article}
\usepackage[margin=1in,top=0.9in,bottom=0.9in]{geometry}
\usepackage{amsmath,amssymb,amsthm,hyperref}
\hypersetup{colorlinks=true,linkcolor=blue,urlcolor=blue,citecolor=blue}
\setlength{\emergencystretch}{2em}
\newtheorem{theorem}{Theorem}
\newtheorem{lemma}[theorem]{Lemma}
\title{An odd-degree identity for polynomial values\\at derivative roots}
\author{Research proof and formalization packet}
\date{October 26, 2026}
\begin{document}
\maketitle
\begin{abstract}
We prove Conjecture~2 of Wojnar, Wojnar and Brin: an alternating binomial
combination of the means of an odd-degree polynomial evaluated at the
roots of its successive derivatives vanishes. The result holds over the
complex numbers, with roots counted with multiplicity. Normalized Newton
moments extend these means polynomially in the remaining derivative
degree. Their artificial zero-degree value is an apolar self-pairing,
which vanishes in odd degree, so a finite difference gives the identity.
The full arbitrary-polynomial statement, including repeated roots and
the source definitions, is verified in Lean. Classical moment,
finite-difference and apolarity ingredients are explicitly credited.
\end{abstract}

\section{The exact identity}
Let $f\in\mathbb C[z]$ have degree $D\geq1$. For $0\leq r<D$ write
\[
 M_r(f)=\frac{1}{D-r}\sum_{f^{(r)}(\zeta)=0}f(\zeta),
\]
where the sum uses the full multiset of $D-r$ roots, including
multiplicity. This is the source mean denoted $\varphi_{D,0,r}$ by
Wojnar--Wojnar--Brin \cite{WWB}.

\begin{theorem}\label{thm:main}
For every complex polynomial of odd degree $D$,
\begin{equation}\label{eq:main}
 \sum_{r=1}^{D-1}(-1)^r\binom Dr M_r(f)=0.
\end{equation}
\end{theorem}

This is Conjecture~2 of the latest listed source revision,
arXiv:1706.08381v2, Section~4, and MathDB~336466. No real-rootedness,
genericity or simple-root assumption is needed. For $D=1$ the sum is
empty. Known individual degrees and the auxiliary identities below are
not separately counted outcomes.

\section{Polynomial continuation of normalized moments}
First assume $f$ is monic. Write
\[
 f(z)=F_D(z)=\sum_{k=0}^D\binom Dk a_k z^{D-k},\qquad a_0=1,
\]
and extend the parameter sequence beyond $D$ arbitrarily. For each
integer $m\geq1$ put
\[
 F_m(z)=\sum_{k=0}^m\binom mk a_k z^{m-k}.
\]
For $1\leq m\leq D$, direct differentiation gives
\begin{equation}\label{eq:appell}
 f^{(D-m)}=\frac{D!}{m!}F_m.
\end{equation}
Thus $F_m$ has exactly the required derivative-root multiset. Write
$S_i(m)$ for the sum of the $i$th powers of its roots.

Define polynomials in an independent variable $T$ by
\[
 B_i(T)=\frac{(T-1)(T-2)\cdots(T-i+1)}{i!},\qquad
 C_i(T)=T B_i(T)\quad(i\geq1),
\]
with empty product one. Set $q_0(T)=1$ and, for $i\geq1$,
\begin{equation}\label{eq:moment}
 q_i(T)=-\sum_{k=1}^{i-1}C_k(T)a_k q_{i-k}(T)-iB_i(T)a_i.
\end{equation}

\begin{lemma}\label{lem:moment}
For every integer $m\geq1$ and $i\geq1$,
\[
 q_i(m)=\frac{S_i(m)}m,\qquad
 \deg q_i\leq i-1,\qquad q_i(0)=(-1)^i a_i.
\]
\end{lemma}
\begin{proof}
The coefficients of $F_m$ are $c_k=\binom mk a_k$, with $c_k=0$ for
$k>m$. Newton's identities, including the homogeneous recurrence beyond
the degree, give for every $i\geq1$
\[
 S_i(m)+\sum_{k=1}^{i-1}c_k S_{i-k}(m)+i c_i=0.
\]
Since $C_i(m)=\binom mi$ and $mB_i(m)=C_i(m)$, division by $m$ and
induction prove the evaluation claim. This also covers $i>m$, which is
essential when evaluating a degree-$D$ polynomial on roots of $F_m$.
In \eqref{eq:moment}, each product $C_kq_{i-k}$ has degree at most
$k+(i-k-1)=i-1$, as does $B_i$. Induction proves the degree bound.
Finally, $C_k(0)=0$ and $B_i(0)=(-1)^{i-1}/i$, yielding the stated value
at zero.
\end{proof}

Now form
\[
 \Phi(T)=\sum_{k=0}^D\binom Dk a_k q_{D-k}(T).
\]
Lemma~\ref{lem:moment}, with $q_0=1$, gives $\deg\Phi\leq D-1$ and
\begin{equation}\label{eq:mean}
 \Phi(m)=\frac1m\sum_{F_m(\zeta)=0}f(\zeta)=M_{D-m}(f)
       \qquad(1\leq m\leq D).
\end{equation}
In particular $\Phi(D)=0$, because $f$ vanishes at every root of $F_D=f$.
At the artificial argument zero,
\begin{equation}\label{eq:apolar}
 \Phi(0)=\sum_{k=0}^D\binom Dk(-1)^{D-k}a_k a_{D-k}.
\end{equation}
This is a value of the polynomial continuation, not a mean over roots
of a constant polynomial. For odd $D$, pairing the $k$ and $D-k$ terms
in \eqref{eq:apolar} shows that $\Phi(0)=0$. This is the classical
odd-degree self-apolarity cancellation.

\section{Finite differences and the source statement}
\begin{proof}[Proof of Theorem~\ref{thm:main}]
The $D$th finite difference of a polynomial of degree at most $D-1$
vanishes, hence
\[
 \sum_{m=0}^D(-1)^m\binom Dm\Phi(m)=0.
\]
Both endpoint terms are zero. Substitute \eqref{eq:mean} in the remaining
sum and reindex $r=D-m$. The resulting sum is
$(-1)^D$ times \eqref{eq:main}; this nonzero common factor can be removed.

For arbitrary $f$, divide by its nonzero leading coefficient $c$ to
obtain a monic polynomial $g$. Every derivative-root multiset is
unchanged by this scaling, and $M_r(f)=cM_r(g)$. Multiplying the monic
identity by $c$ proves the full statement.
\end{proof}

\section{Classical ingredients, novelty, and formal scope}
The normalized root-moment and finite-difference phenomenon has important
older precedents. Toda \cite{Toda36} studies finite differences of
test-function means over successive derivative roots; his pages~315--316
use binomial-normalized coefficients and Waring's formula. Popoviciu
\cite{Pop} explicitly records Toda's vanishing result for a test
polynomial of degree at most $n$, with $n$ strictly smaller than the
degree of the root polynomial. Those are close precedents for the
moment-continuation ingredient and receive credit here.

The present target has the same-degree self-test and an omitted
zero-degree root mean. The needed boundary value is supplied by the
apolar continuation \eqref{eq:apolar}; the inspected older statements do
not explicitly give this endpoint. Odd self-apolarity itself is also
classical, for example Ottaviani \cite{Ott}, Section~4.3, Proposition~18.
The proposed contribution is the complete source-mean bridge and
resulting all-degree resolution of \cite{WWB}, not any of these classical
tools or an exhaustive classification of all mean relations.

The detailed primary-source investigation found no previous full proof
of the exact source conjecture. Toda's 1934 precursor \cite{Toda34}
remained inaccessible in full because its public publisher endpoint
returned an anti-bot page. Its 1936 successor and Popoviciu's original
1940 discussion were read at the formula level, including rendered
pages, but do not exclude every possible additional formula in the
unread precursor. This material retrieval limitation is explicit in
the independent novelty review. Acceptance under the documented
no-known-prior-resolution standard is not a certification of worldwide
priority.

The six-module Lean certificate inventories 50 declarations. It proves
Newton's identities on complex root multisets, all moment continuation
claims including indices above the remaining degree, exact iterated
derivative/root correspondence, endpoint cancellation, finite
differences, interval trimming and arbitrary-polynomial normalization.
The final theorem
\texttt{OddRootMeans.source\_conjecture\_336466} assumes only that the
complex polynomial's natural degree is odd. No user-supplied moment,
splitting, source-bridge or distinct-root hypothesis remains. Fresh
independent replays use pinned Lean~4.34.1 and only \texttt{propext},
\texttt{Classical.choice}, and \texttt{Quot.sound}; no proof holes or
additional axioms occur. Numerical trace experiments are supplementary
and not dependencies of the proof.

\clearpage
\begin{thebibliography}{9}
\bibitem{WWB} G. G. Wojnar, D. Sz. Wojnar and L. Q. Brin,
\emph{Universal Linear Mean Relationships in All Polynomials},
\href{https://arxiv.org/abs/1706.08381v2}{arXiv:1706.08381v2} (2017),
Section~4, Conjecture~2; root conventions in Section~1.
\bibitem{Toda36} K. Toda, \emph{A Method of Approximation of Convex Functions},
Tohoku Math. J. 42 (1936), 311--317.
\href{https://www.jstage.jst.go.jp/article/tmj1911/42/0/42_0_311/_pdf/-char/en}{Primary text}.
\bibitem{Pop} T. Popoviciu, \emph{Notes sur les fonctions convexes d'ordre
sup\'erieur (III)}, Mathematica 16 (1940), 74--86, Section~11.
\href{https://ictp.acad.ro/notes-on-higher-order-convex-functions-iii/}{Institute-hosted original}.
\bibitem{Ott} G. Ottaviani, \emph{Five Lectures on Projective Invariants},
\href{https://arxiv.org/abs/1305.2749}{arXiv:1305.2749} (2013),
Section~4.3, Proposition~18.
\bibitem{Toda34} K. Toda, \emph{On Certain Functional Inequalities},
J. Sci. Hiroshima Univ., Ser. A 4 (1934), 27--40.
\href{https://doi.org/10.32917/hmj/1558749761}{DOI: 10.32917/hmj/1558749761}.
Full body not inspected; see the stated novelty limitation.
\end{thebibliography}
\end{document}
